\section{Measurement Protocol}
\label{sec:method}

\subsection{Scenario Sets}
All sets use seed 11 unless stated otherwise, and numbers from different sets are never mixed in one comparison.
\begin{itemize}
    \item \emph{Aging set} (Tables~\ref{tab:aging}--\ref{tab:tracking}, Figs.~\ref{fig:aging}--\ref{fig:scaling}): 400 scenarios with walking users and 20 point scatterers, close to the 15--19 clusters of the 3GPP indoor-office model; with $N=64$, the same construction with an $8\times 8$ panel.
    \item \emph{Beam set} (Section~\ref{sec:res_beam}): 400 more scenarios, 6 scatterers.
    \item \emph{Static set} (Fig.~\ref{fig:eqms}): 6000 static scenarios, randomly split 4200/900/900 (training/validation/test); the learned update is validated on 500, Adam on 400 of these (on all 500 it picks the same setting); results use 600 test scenarios.
\end{itemize}
The generator is not prefix-stable; trained models are evaluated at their training sample count.

\subsection{Methods}
\begin{table}[t]
\caption{Methods, settings and computation cost (aging set unless stated).}
\label{tab:hyper}
\footnotesize
\setlength{\tabcolsep}{3pt}
\begin{tabular}{@{}>{\raggedright\arraybackslash}p{0.2\columnwidth}>{\raggedright\arraybackslash}p{0.51\columnwidth}>{\raggedright\arraybackslash}p{0.23\columnwidth}@{}}
\toprule
Method & Settings & Cost \\
\midrule
SU-average & phase of the sum of the $K$ single-user alignments & 0.42~ms \\
SU-best-of & best of the $K$ alignments and SU-average & 2.5~ms \\
Adam & start SU-average; soft-min, $\beta=8$; step size 0.5 (static set: 0.75, selected for $L\le 16$); $L^\star=4$--6 at 1.5~m/s & 1.66~ms/step$^{a}$ \\
Armijo & start SU-average; soft-min; initial step 3, backtracking; $L^\star=4$ & 2.09~ms/step \\
Coordinate descent & 16 phase levels, hard min, 1--8 sweeps (best: 2) & 107~ms/sweep \\
Long-run Adam & Adam, step size 0.05, 1000 iterations, 3 random starts, hard min & 5.85~s \\
Learned update (static set) & $L\in\{1,2,4,8,16\}$, one network per $L$; 41 input channels; 272{,}835--272{,}850 parameters; unsupervised soft-min training (Adam, $10^{-3}$, cosine schedule, batch 32, 12 epochs) & 2.7--4.0 Adam steps per step \\
\bottomrule
\end{tabular}

\vspace{2pt}\raggedright\footnotesize $^{a}$1.66~ms in Sections~\ref{sec:res_aging}--\ref{sec:res_tracking}; 1.60--1.71~ms (static set) in Section~\ref{sec:res_reversal}; 1.33~ms (re-measurement) in Section~\ref{sec:res_robust}; other sessions: supplementary material.
\end{table}
\emph{SU-average} and \emph{SU-best-of} are closed-form single-user phase alignments; the latter keeps the best of $K{+}1$ candidates. \emph{Adam}~\cite{Kingma2015ICLR} optimizes the soft-min surrogate, starting from SU-average. \emph{Armijo} is steepest ascent with Armijo's rule~\cite{Armijo1966PJM} on the same surrogate, with a per-scenario backtracking line search. \emph{Coordinate descent} searches each element over 16 phase levels, for up to 8 sweeps; \emph{long-run Adam} runs 1000 iterations from 3 random starts. The \emph{learned update} is a deep-unfolded network with $L$ learned steps; one network is trained per $L$, without labels. For each user and element it receives 37 static input channels: the real and imaginary parts of the normalized RIS--user channel (2), of the BS--RIS channel at each of the $N_t$ antennas (16) and of the direct channel (16), the cosine and sine of the user's single-user alignment phase, and the logarithm of the user's relative channel power. At every step four dynamic channels are added: the cosine and sine of the current phase, the user's standardized rate and the normalized gradient of the soft-min rate. A convolutional and long short-term memory (LSTM) backbone along the elements, two layers that exchange information across users and a per-element head produce a phase update, scaled by a learned step size; the weights are shared across steps. Table~\ref{tab:hyper} lists the settings of all methods.

\emph{Tuning.} On the static set, Adam's step size and objective (soft-min or hard min) are selected on the validation set at every budget: from ten values in $[0.02, 2.0]$ for $L\le 16$, where 0.75 with the soft-min, an interior point, is selected every time, and from six values in $[0.2, 1.5]$ for $L=24$--$64$. On the aging set Adam uses step size 0.5 (0.25, 0.75 and 1.0 give lower peaks), and Armijo's initial step is chosen from five values (3 is selected); both choices use the reported scenarios, so they are robustness checks rather than held-out selections.

\subsection{Timing}
Latencies are wall-clock times for one scenario on an Intel Xeon CPU at 2.1~GHz with one PyTorch thread (except the two-thread 1.66~ms below), in unoptimized Python (Section~\ref{sec:res_scaling} considers a faster Adam step); the median over repetitions is kept, and most measurements are repeated in several rounds. Adam's step time appears with three values, each tied to one measurement session (Table~\ref{tab:hyper}). The aging-set results of Sections~\ref{sec:res_aging}--\ref{sec:res_tracking} use 1.66~ms (13.3~ms for 8 steps), from an earlier two-thread measurement made before a precoder fix; it lies within the 1.49--2.02~ms of later rounds on the aging set. The static-set comparison (Section~\ref{sec:res_reversal}) uses its own five-round medians, 1.60--1.71~ms per step. Section~\ref{sec:res_robust} uses a one-thread re-measurement of 1.33~ms (load average 0.45 before timing; coordinate descent 0.70~s); with it $L^\star=6$, $L=48$ costs 14.9\%, and SU-best-of stays ahead of 8-sweep coordinate descent by 0.26~b/s/Hz (CI 0.15--0.37; optimistic for coordinate descent), so the optimum and the cost of $L=48$ do not depend on which of the two aging-set values is used. Armijo costs 1.26 times as much as Adam per step (2.09~ms; per-scenario ratio 0.81--1.97). On the static set the learned network's input features are excluded from its time, which favors it. Other checks use separate sessions; load averages, round counts and spreads are in the supplementary material.

\subsection{Statistics and Pilots}
Effective rates are interpolated on a time grid at each method's latency plus an optional acquisition delay $T_{\mathrm{ov}}$. Evaluating at the exact latency instead changes Table~\ref{tab:aging} by at most 0.006~b/s/Hz within the reliable horizon (0.008 and 0.037 beyond it) and the Adam sweep of Fig.~\ref{fig:optimum} by at most 0.027; the optimum stays at 4--6 steps and the cost of $L=48$ becomes 14.2\%. Because absolute rates vary strongly across scenarios ($\pm$0.17--0.24~b/s/Hz for the means of Table~\ref{tab:aging}), methods are compared with paired bootstrap confidence intervals (CIs) over scenarios (2000 resamples). Differences of a few thousandths of a b/s/Hz between neighboring $L$ are not resolved. Best configurations are selected and reported on the same scenarios unless stated, without a multiple-comparison correction, and should be read as estimates.

We tie $T_{\mathrm{ov}}$ and the pilot time $T_{\mathrm{p}}$ to cascaded channel estimation~\cite{Zheng2022COMST}. With the uplink scheme of~\cite{Wang2020TWC} and TDD reciprocity, 48 pilot symbols suffice for $K=4$, $N=32$ and $N_t=8$ in the noise-free case. With 5G New Radio (NR) numerology~\cite{3gpp_nr_phy} and one orthogonal frequency-division multiplexing (OFDM) symbol per pilot they last 3.43, 1.71 and 0.86~ms at 15, 30 and 60~kHz subcarrier spacing; feedback, estimation processing and RIS switching come on top and are not modeled. The periodic controller is simulated on a 0.5~ms grid over 200~ms, with update periods of 10--80~ms or back-to-back, averaging over 100--200~ms; pilot airtime is charged as the fraction of each grid step it occupies, each phase takes effect at its exact deployment instant, and channel snapshots are taken at the nearest grid point.

\emph{Reproducibility.} One script generates every table and figure from stored result files, except Table~\ref{tab:tracking} (a corrected copy of the tracking simulation) and Section~\ref{sec:res_robust} (separate scripts); rerunning the aging sweeps reproduces them bit for bit. A separate audit script checks 126 assertions about the learning experiments of the underlying thesis; among the results here it covers those of Section~\ref{sec:res_sensitivity}.
